%! Skills Focus: Selecting an Appropriate Inference Procedure for Categorical Data — Unit 8, Lesson 8.7 — Group Activity (Answer Key)
\documentclass[10pt]{article}
\usepackage{apstats-article}
\usepackage{apstats-key}   % pulls in -boxes; do NOT also load -boxes

\begin{document}

\pageheader{Unit 8, Lesson 8.7}{Group Activity --- Answer Key}
\namepartnerperiod

\vspace{0.08in}

\textbf{Procedure Menu} (use for Tier R): one-sample $z$-interval for $p$ \;|\;
one-sample $z$-test for $p$ \;|\; two-sample $z$-test for $p_1-p_2$ \;|\;
chi-square goodness-of-fit test \;|\;
chi-square test for homogeneity \;|\;
chi-square test for independence

\vspace{0.06in}

\begin{scenariobox}[Scenario 1 --- Streaming Preferences]{sky}
A media company takes one random sample of 400 adults and records two things: each person's
age group (Under 35 or 35 and Over) and their primary streaming service (Service A, B, or C).
The company wants to test whether there is an association between age group and streaming
service preference.

\textbf{Tier R:} Name the procedure and state one key condition.
\begin{itemize}
  \item Procedure: \ans{chi-square test for independence}
  \item Key condition: \ans{All expected cell counts must be at least 5.}
\end{itemize}

\textbf{Tier A:}
\begin{enumerate}
  \item Procedure name: \ans{chi-square test for independence}
  \item $H_0$: \ansline{There is no association between age group and primary streaming service preference.}
  \item $H_a$: \ansline{There is an association between age group and primary streaming service preference.}
  \item Key condition to check: \ansline{Expected counts condition: all expected cell counts $\ge 5$. (Cannot verify without observed counts, but $n = 400$ with a $2 \times 3$ table makes this very likely.)}
\end{enumerate}
\end{scenariobox}

\begin{scenariobox}[Scenario 2 --- M\&M Color Distribution]{sky}
A statistics student purchases a large bag of M\&Ms. The manufacturer claims the color
distribution is 24\% blue, 20\% orange, 16\% green, 14\% yellow, 13\% red, and 13\% brown.
The student counts colors in a random sample of 200 M\&Ms from the bag and wants to test
whether her bag matches the claimed distribution.

\textbf{Tier R:} Name the procedure and state one key condition.
\begin{itemize}
  \item Procedure: \ans{chi-square goodness-of-fit test}
  \item Key condition: \ans{All expected counts must be at least 5.}
\end{itemize}

\textbf{Tier A:}
\begin{enumerate}
  \item Procedure name: \ans{chi-square goodness-of-fit test}
  \item $H_0$: \ansline{The distribution of M\&M colors is 24\% blue, 20\% orange, 16\% green, 14\% yellow, 13\% red, and 13\% brown.}
  \item $H_a$: \ansline{The distribution of M\&M colors differs from the claimed distribution.}
  \item Minimum expected count for blue: $200 \times 0.24 =$ \ans{48} \quad Condition met? \ans{yes ($48 \ge 5$)}
\end{enumerate}
\end{scenariobox}

\begin{scenariobox}[Scenario 3 --- Vaccination Status and Opinion]{sky}
A public health researcher randomly samples 150 vaccinated adults and 150 unvaccinated adults.
She asks each person whether they support a new public health mandate (Support, Neutral, or
Oppose). She wants to test whether the distribution of opinion is the same for vaccinated and
unvaccinated adults.

\textbf{Tier R:} Name the procedure.
\begin{itemize}
  \item Procedure: \ans{chi-square test for homogeneity}
\end{itemize}

\textbf{Tier A:}
\begin{enumerate}
  \item Procedure name: \ans{chi-square test for homogeneity}
  \item $H_0$: \ansline{The distribution of opinion (Support, Neutral, Oppose) is the same for vaccinated and unvaccinated adults.}
  \item $H_a$: \ansline{The distribution of opinion differs between vaccinated and unvaccinated adults.}
  \item Explain in one sentence why this is \emph{not} a test for independence: \ansline{The data come from two separate random samples (one vaccinated, one unvaccinated), not from a single random sample with two variables recorded.}
\end{enumerate}
\end{scenariobox}

\begin{scenariobox}[Scenario 4 --- Voter Turnout]{sky}
A political researcher randomly selects 500 registered voters and records whether each voted
in the last election (yes or no). He wants to test whether the proportion who voted is
greater than 55\%.

\textbf{Tier A:}
\begin{enumerate}
  \item Procedure name: \ans{one-sample $z$-test for $p$}
  \item $H_0$: \ans{$p = 0.55$} \quad $H_a$: \ans{$p > 0.55$}
  \item Large Counts condition: $500 \times 0.55 =$ \ans{275} $\ge 10$? \ans{yes} \quad
        $500 \times 0.45 =$ \ans{225} $\ge 10$? \ans{yes}
\end{enumerate}
\end{scenariobox}

\begin{scenariobox}[Scenario 5 --- Flu Shot Rates]{sky}
A hospital system randomly samples 200 patients from its urban clinic and 200 patients from
its rural clinic. It records whether each patient received a flu shot this season (yes or no).
The hospital wants to test whether the proportion who received a flu shot differs between the
two clinics.

\textbf{Tier A:}
\begin{enumerate}
  \item Procedure name: \ans{two-sample $z$-test for $p_1 - p_2$}
  \item $H_0$: \ans{$p_1 = p_2$} \quad $H_a$: \ans{$p_1 \ne p_2$}
  \item Why is this a two-sample $z$-test rather than a chi-square test for homogeneity?
        \ansline{The outcome is binary (yes/no flu shot), so there are only two categories. A two-sample $z$-test is the appropriate procedure when comparing two proportions; chi-square homogeneity is used when the response has three or more categories.}
\end{enumerate}
\end{scenariobox}

\begin{scenariobox}[Scenario 6 --- Tier E: Compare and Contrast]{sky}
\textbf{Tier E only.}

A researcher takes a random sample of 300 college students and records two variables: class
year (First-Year, Sophomore, Junior, Senior) and whether each student has a declared major
(yes or no).

\begin{enumerate}
  \item A student says this situation calls for a two-sample $z$-test for $p_1 - p_2$
        because the outcome (declared major) is binary. Identify the error in this reasoning
        and name the correct procedure.
        \ansline{Error: the two-sample $z$-test requires exactly two groups (samples). Here, class year has four categories, not two, and there is only one sample.}
        \ansline{Correct procedure: chi-square test for independence (one sample, two variables recorded).}
  \item A second student says this calls for a chi-square goodness-of-fit test. Identify the
        error and explain why.
        \ansline{Error: a GOF test compares a single variable's observed distribution to a claimed distribution. Here, there is no claimed distribution, and we have two variables.}
        \ansline{The correct procedure is chi-square independence because we are examining the association between two categorical variables.}
  \item State the correct $H_0$ and $H_a$ for the appropriate procedure.
        \begin{itemize}
          \item $H_0$: \ansline{There is no association between class year and whether a student has a declared major.}
          \item $H_a$: \ansline{There is an association between class year and whether a student has a declared major.}
        \end{itemize}
\end{enumerate}
\end{scenariobox}

\end{document}
